\documentclass[12pt]{article}
\usepackage[utf8]{inputenc}
\usepackage[T1]{fontenc}
\usepackage[margin=1in]{geometry}
\usepackage{setspace}
\usepackage{xcolor}
\usepackage[colorlinks=true, urlcolor=blue!60!black, linkcolor=blue!60!black]{hyperref}
\setlength{\parindent}{0em}
\setlength{\parskip}{0.6em}
\renewcommand{\baselinestretch}{1.0}
\begin{document}
\begin{center}
{\LARGE \textit{Pierce v. Society of Sisters}, 268 U.S. 510 (1925)\par}
\vspace{0.5em}
{\large Case Brief --- Two-Page Edition\par}
\end{center}

\section{Memory Jogger}
Oregon could not require all children to attend public schools: parents have a fundamental liberty to direct their children's education, and the state may not ``standardize its children'' (\href{https://supreme.justia.com/cases/federal/us/268/510/}{Justia}).

\section{Detailed Case Facts}
Amid post-war nativism and anti-Catholic sentiment, Oregon voters in 1922 required every child aged eight to sixteen to attend public school, effectively outlawing private and parochial primary schools. A Catholic teaching order, the Society of Sisters, and Hill Military Academy, a secular for-profit school, sued as parents began withdrawing students before the law took effect.

\section{Procedural History}
A three-judge district court enjoined enforcement, holding that the schools' property was protected and their suits ripe. Oregon appealed directly.

\section{Judicial Votes}
Affirmed, 9--0. Justice McReynolds wrote for a unanimous Court.

\section{Holding}
Under \textit{Meyer v. Nebraska}, compelling public-school attendance unreasonably interferes with parents' liberty to direct their children's upbringing and education. Private schools may sue before the law's effective date to protect business and property threatened by compulsion of their patrons.

\section{Analysis of Opinions}
\textbf{Opinion of the Court:} The state may require attendance at some school and reasonably regulate all schools---a concession that is technically dicta---but it may not destroy useful private schools when no emergency requires it. The law bore no reasonable relation to a legitimate state purpose. Though corporations lack Fourteenth Amendment liberty, the schools' ``business and property'' were threatened, the law was no mere amendment of corporate charters, and present losses made the suits ripe. No Justice wrote separately.

\section{Examples: Future Applications}
\textbf{Same-Side:} a ban on private middle schools to promote civic cohesion is invalid. \textbf{Opposite-Side:} fire codes, minimum instructional days, and core-subject requirements for private schools are valid regulation. \textbf{Fence-Sitter:} a substantial-equivalency mandate that would close yeshivas centered on Talmudic study.

\section{Critique}
Originalists note its \textit{Lochner}-era roots; Justice Scalia called it a relic of repudiated substantive due process, and others would ground it in the First Amendment. Child-rights scholars such as Woodhouse charge that it treats the child as parental property.

\section{Key Quotations}
\begin{quote}
``Under the doctrine of \textit{Meyer v. Nebraska}, 262 U.\,S. 390, we think it entirely plain that the Act of 1922 unreasonably interferes with the liberty of parents and guardians to direct the upbringing and education of children under their control. As often heretofore pointed out, rights guaranteed by the Constitution may not be abridged by legislation which has no reasonable relation to some purpose within the competency of the State.''
\end{quote}
\noindent --- Justice McReynolds, for the Court, 268 U.S. at 534--35.

\end{document}
